\section{The Italian Decade: Orvieto and Rome}

\subsection{Lector at Orvieto; the Catena aurea}

From 1259/61 the order stationed Thomas in its Roman province---the
sources name Anagni, Orvieto, Rome, Viterbo without a fully secure
year-by-year map---and at Orvieto, where Urban~IV held his court, he
served as conventual lector, teacher of the ordinary friars. The papal
court put him to work. For Urban he began the \work{Catena aurea}, the
``golden chain'': a continuous gloss on the four Gospels woven entirely
from the Fathers, Greek as well as Latin, its Matthew portion dedicated
to Urban himself (so the 1912 Encyclopedia, which cites the work as proof
of his ``prodigious memory''---a judgment about the sources he commanded,
since the chain required hunting Greek patristic material then almost
unknown in the West). At the same court he wrote \work{Contra errores
Graecorum} on the disputed questions between the Latin and Greek
churches---a work honest history must qualify: it depends on a dossier of
patristic quotations supplied to him, some of them spurious, a limit of
his sources rather than his honesty.

To Orvieto in 1264 belongs also the act that tied Thomas's name forever
to the Eucharist: Urban~IV's bull \work{Transiturus} instituted the feast
of Corpus Christi for the whole Church, and tradition holds that the pope
asked Thomas to compose its Mass and office---the texts from which come
\textit{Pange lingua}, \textit{Lauda Sion}, and the rest of the Church's
classic Eucharistic poetry.

\witnesslimit{The attribution deserves exactness. The papal-commission
tradition is canonization-era; by 1912 the Encyclopedia could report it
as ``established beyond doubt'' on Mandonnet's authority; twentieth-century
liturgical scholarship then complicated the question---Lambot's study of
the office's origins (1942) and Gy's re-examination (1982) frame the
modern debate over how the Roman office relates to an earlier Li\`ege
office and what precisely Thomas composed or reworked. Those studies are
identified here at bibliographic level only, so this study asserts the
early tradition and the debate, not an item-by-item authorship. The
Church's own practice is instructively careful: when the Catechism quotes
the hymn \textit{Adoro te devote}, its footnote attributes it to
``St.~Thomas Aquinas (attr.)''---attributed---and no more (CCC 1381
note).}

\subsection{The office of the wise man}

To this Italian period belongs the completion of the \work{Summa contra
Gentiles}, the four-book exposition of Christian truth argued, wherever
possible, on grounds a philosophical opponent could share. Its preface
contains the nearest thing Thomas ever wrote to a personal manifesto.
Taking up Hilary's words, he makes them his own: aware that he owes his
life's ``chief office'' to God, he declares---\latin{``Propositum nostrae
intentionis est veritatem quam fides Catholica profitetur, pro nostro
modulo manifestare, errores eliminando contrarios''}: the purpose of our
intention is to make plain, according to our small measure, the truth
that the Catholic faith professes, eliminating the contrary errors; and
so, ``taking confidence from God's mercy to pursue the office of the wise
man, though it exceeds our own powers'' (\latin{``Assumpta igitur ex
divina pietate fiducia sapientis officium prosequendi, quamvis proprias
vires excedat\ldots''}), he begins (\work{Summa contra Gentiles} I.2,
Leonine text as presented by the Corpus Thomisticum). It is the whole
man in two sentences: the work is an office, the measure is small, the
confidence is borrowed.

\subsection{Rome: the Summa begun}

About 1265 the order set Thomas over its house of studies at Santa Sabina
in Rome, with young Dominicans to form---and it was for them, not for the
schools' virtuosi, that he conceived his masterwork. The \work{Summa
theologiae} announces itself in its prologue as a book for beginners: an
ordering of the whole of sacred doctrine, briefly and clearly, ``as the
matter will allow,'' stripped of the useless questions, digressions, and
repetitions that (he says) weary students in the existing books. The plan
is a single arc, stated at the head of the second question in words
Benedict~XVI's catechesis on the Summa quotes: ``we shall treat: (1) Of
God; (2) Of the rational creature's advance towards God; (3) Of Christ,
Who as man, is our way to God''---all things from God, all things back to
God, through Christ.

The first article states the founding claim of the whole, in the sentence
this study reproduces from the Leonine edition: \latin{``Respondeo
dicendum quod necessarium fuit ad humanam salutem, esse doctrinam quandam
secundum revelationem divinam, praeter philosophicas disciplinas, quae
ratione humana investigantur''}---``It was necessary for man's salvation
that there should be a knowledge revealed by God besides philosophical
science built up by human reason'' (\work{Summa theologiae} I.1.1,
Leonine ed., vol.~4, 1888; English Dominican translation). Reason keeps
its whole domain; revelation is necessary because man is ordered to an
end that exceeds it. On that double affirmation---the integrity of nature
and reason, the primacy and gratuity of grace and revelation---the entire
structure stands, and Benedict~XVI's catecheses present exactly this
mutual friendship of faith and reason as Thomas's permanent achievement.
The First Part was written in Italy; the immense Second Part, on human
action, virtue, law, and grace, belongs mostly to the second Paris
regency; the Third, on Christ and the sacraments, to Paris and Naples;
and the whole would never be finished.
